\documentclass[10pt,a4paper]{amsart}
\usepackage[margin=19mm]{geometry}
\usepackage{amsmath,amssymb,mathtools}
\usepackage[scr=rsfs,cal=euler]{mathalfa}
\usepackage{xfrac}
\usepackage[unicode,hidelinks,bookmarks=false]{hyperref}
\newcommand{\ie}{i.e.\ }
\newcommand{\cf}{cf.\ }
\newcommand{\resp}{respectively\ }
\newcommand{\Subsection}{\subsection}
\providecommand{\nobreakdash}{\nobreak\hbox{-}}
\setlength{\emergencystretch}{2em}
\multlinegap0pt
\usepackage{natbib}
\usepackage{amssymb}
\usepackage[shortlabels]{enumitem}
\usepackage{bm}
\usepackage{algorithm}\usepackage{algpseudocode}
\usepackage{booktabs}
\usepackage{xcolor}
\usepackage{caption}
\newtheorem{assumption}{Assumption}
\newtheorem{lemma}{Lemma}
\newtheorem{theorem}{Theorem}
\newcommand{\ebb}{\mathbb{E}}  %
\newcommand{\rbb}{\mathbb{R}}  %
\title{Generalized Stochastic Gradient Descent with Momentum Methods for Smooth Optimization}
\author{Zimeng Wang, Alp Yurtsever}
\date{Source: arXiv:2602.23444v1 (CC BY 4.0)}
\begin{document}
\maketitle

\noindent\emph{Source and licence.} Generalized Stochastic Gradient Descent with Momentum Methods for Smooth Optimization. Version arXiv:2602.23444v1 (CC BY 4.0), distributed by arXiv under the Creative Commons Attribution 4.0 International licence, http://creativecommons.org/licenses/by/4.0/, as recorded in the arXiv licence metadata for this exact version and shown on the abstract page https://arxiv.org/abs/2602.23444v1. Full licence notice: \texttt{licenses/arxiv-2602.23444-CC-BY-4.0.txt}.\par\smallskip

\noindent\textit{Editorial scope.} Counted unit: the article's theorem thm:pl (source0.tex lines 946--956). The article's complete original proof of it is delivered with it (source0.tex lines 958--1061, one \texttt{proof} environment of 104 lines). No mathematics is written, repaired or re-derived: the statement and the proof are the article's own lines, in the article's own notation, and no retained line is substituted or re-flowed.  Retained uncounted as the source-local context the proof needs: lines 172--178; lines 279--284; lines 296--298; lines 799--801; lines 803--808; lines 875--882; lines 935--937; lines 941--943.  Nothing was omitted from any retained span, so the package carries no omission record.  No retained line contains a citation, so the wrapper carries no bibliography and no bibliography entry.  No retained line contains a figure environment, an image-inclusion command or drawing code, so no image is delivered and the package declares no asset dependency.  Editorial formatting only, in addition: an AMS \texttt{amsart} preamble that restates the article's own declarations and macros used by the retained lines, namely the environments \texttt{assumption}, \texttt{lemma}, \texttt{theorem} on separate counters, together with the standard packages those lines need.  The retained environments print in this document as Assumption 1, Lemma 1, Theorem 1 in order of appearance, whereas the article prints the same items with the numbering of its own full document.  The complete native source of the article as distributed in the arXiv e-print for this exact version is bundled unchanged as \texttt{sources/195401/source0.tex}.
\begin{equation}\label{alg:g-sgdm}
\tag{G-SGDM}
    \begin{aligned}
        &m_k=\beta_k m_{k-1}+(1-\beta_k)g_k, \\
        &x_{k+1}=x_k-\gamma_k g_k-\eta_k m_k,
    \end{aligned}
\end{equation}

\begin{assumption}\label{ass:s-bv}
    The objective function $f$ is $L$-smooth. Moreover, the gradient approximations $\{g_k\}$ are unbiased and have uniformly bounded variance $\sigma^2\ge 0$, i.e.,
    \[
    \ebb_k[g_k]=\nabla f(x_k),\quad \ebb_k[\|g_k-\nabla f(x_k)\|^2]\leq \sigma^2,~\forall k\ge 1.
    \]
\end{assumption}

\begin{equation}\label{def:wk}
w_k= \begin{cases}x_1, \quad \text{if } k=1 &  \\ \frac{1}{1-\beta} (x_k-\beta x_{k-1}+\beta \gamma g_{k-1}), & \text{if } k \ge 2.\end{cases}
\end{equation}

\begin{equation}\label{eq:wk-xk}
    w_k=x_k-\frac{\beta\eta}{1-\beta} m_{k-1},~\forall k\ge 1.
\end{equation}

\begin{lemma}\label{lem:m-var}
Let Assumption \ref{ass:s-bv} hold. Let $\{x_k\}_{k\ge 1}$ and $\{m_k\}_{k\ge 0}$ be the sequences generated by \eqref{alg:g-sgdm} with constant momentum parameter $\beta\in [0,1)$. Then for any $k\ge 0$, it holds that
\[
\ebb[\|m_k\|^2] \leq 2(1-\beta) \sum_{j=1}^k \beta^{k-j}\ebb[\|\nabla f(x_j)\|^2]+\frac{2(1-\beta)}{1+\beta} \sigma^2.
\]
\end{lemma}

\begin{equation}\label{eq:nc-3}
\begin{aligned}
\ebb[f(w_{k+1})] \leq & \ebb[f(w_k)]+\frac{\rho L^2\beta^2\eta^2(\gamma+\eta)}{2(1-\beta)^2}  \ebb[\|m_{k-1}\|^2] \\
& +\Big(\frac{1}{2 \rho}-1\Big)(\gamma+\eta) \ebb[\|\nabla f(x_k)\|^2]+\frac{L (\gamma+\eta)^2}{2} \ebb[\|g_k\|^2]\\
\leq & \ebb[f(w_k)]+\frac{\rho L^2\beta^2\eta^2(\gamma+\eta)}{2(1-\beta)^2} \ebb[\|m_{k-1}\|^2]+\frac{L(\gamma+\eta)^2}{2} \sigma^2 \\
& +\Big(\frac{1}{2 \rho}-1+\frac{L(\gamma+\eta)}{2}\Big)(\gamma+\eta) \ebb[\|\nabla f(x_k)\|^2].
\end{aligned}
\end{equation}

\begin{equation}\label{eq:pl-condition}
\| \nabla f(x) \|^2 \ge 2 \mu(f(x)-f^*),~\forall x\in\rbb^d.
\end{equation}

\begin{equation}\label{def:varphi}
\varphi_k=f(w_k)-f^*+C \sum_{j=1}^{k-1} \beta^{k-1-j}\|\nabla f(x_j)\|^2,
\end{equation}

\begin{theorem}\label{thm:pl}
Let Assumption \ref{ass:s-bv} hold and suppose that $f$ satisfies the $\mu$-PL condition \eqref{eq:pl-condition}. Let $\{x_k\}_{k\ge 1}$ be the sequence generated by \eqref{alg:g-sgdm} with constant parameters $\beta\in [0,1)$, $\gamma\ge 0$, and $\eta>0$. If we choose
\begin{equation}\label{eq:pl-lr}
    0<\gamma+\eta \leq \min\Big\{\frac{1-\beta}{2 L},~\frac{1-\beta}{8 \beta^2 L}\Big\},
\end{equation}
then for any $t\ge 1$, it holds that
\[
\ebb[f(w_{t+1})-f^*]\leq\Big(1-\frac{\mu(\gamma+\eta)}{18}\Big)^{t}(f(x_1)-f^*)+\Big(\frac{3 \beta^2 \eta}{1+\beta}+\gamma+\eta\Big) \frac{9L}{\mu} \sigma^2,
\]
where $w_t$ is defined in \eqref{def:wk}.
\end{theorem}

\begin{proof}[Proof of Theorem \ref{thm:pl}]
For notational simplicity, denote
\[
    M:=\gamma+\eta-\frac{L(\gamma+\eta)^2}{1-\beta}-\frac{L(\gamma+\eta)^2}{2}.
\]
Consider the Lyapunov functions $\{\varphi_k\}_{k\ge 1}$ defined by \eqref{def:varphi} with
\begin{equation}\label{def:C}
C:=\frac{\frac{\beta^2 L \eta^2}{2}+\frac{M \mu(\mu+L) \beta^2 \eta^2}{1-\beta}}{1-\beta-\mu M+\frac{\mu(\mu+L) \beta^2 \eta^2}{1-\beta}}.
\end{equation}
We first establish some auxiliary inequalities regarding the introduced constants $C$ and $M$. To start off, it follows from $0<\gamma+\eta \leq \frac{1-\beta}{2 L}$ and $\beta\in[0,1)$ that
\begin{equation}\label{eq:pl-M>}
    M=(\gamma+\eta)\Big(1-\frac{(3-\beta) L}{2(1-\beta)}(\gamma+\eta)\Big) \ge \Big(1-\frac{3-\beta}{4}\Big)(\gamma+\eta) \ge \frac{1}{4}(\gamma+\eta)>0.
\end{equation}
On the other hand, we note from $\gamma+\eta\leq\frac{1-\beta}{2 L}$ and $\mu\leq L$ that
\begin{equation}\label{eq:pl-M<}
    M \leq \gamma+\eta \leq \frac{1-\beta}{2 L}\leq\frac{1-\beta}{2 \mu}.
\end{equation}
Hence we derive $C>0$ by noting from \eqref{eq:pl-M<} that
\begin{equation}\label{eq:pl-M-1}
1-\beta-\mu M \ge \frac{1-\beta}{2}\ge 0.
\end{equation}
Next, we show that $M-C$ is positive and bounded from both sides. To this end, first note from $0\leq\eta\leq\gamma+\eta \leq \frac{1-\beta}{8 \beta^2 L}$, \eqref{eq:pl-M-1}, and \eqref{eq:pl-M>} that
\begin{equation}\label{eq:pl-M-C-1}
\begin{aligned}
M(1-\beta-\mu M)-\frac{\beta^2 L\eta^2}{2} \ge& \frac{1-\beta}{2} M-\frac{\beta^2 L \eta^2}{2} \ge \frac{1-\beta}{8}(\gamma+\eta)-\frac{\beta^2 L}{2}(\gamma+\eta)^2 \\
= & (\gamma+\eta)\Big(\frac{1-\beta}{8}-\frac{\beta^2 L}{2}(\gamma+\eta)\Big) \ge \frac{1-\beta}{16}(\gamma+\eta).
\end{aligned}
\end{equation}
On the other hand, we obtain from $\mu\leq L$ and \eqref{eq:pl-lr} that
\begin{equation}\label{eq:pl-M-C-2}
\begin{aligned}
1-\beta-\mu M+\frac{\mu(\mu+L) \beta^2 \eta^2}{1-\beta} & \leq 1-\beta+\frac{\mu(\mu+L) \beta^2(\gamma+\eta)^2}{1-\beta} \\
& \leq 1-\beta+\frac{2L^2 \beta^2}{1-\beta}\frac{1-\beta}{2 L}\frac{1-\beta}{8 \beta^2 L} \\
& = 1-\beta+\frac{1-\beta}{8} = \frac{9}{8}(1-\beta).
\end{aligned}
\end{equation}
Then it follows from $C>0$, \eqref{eq:pl-M<}, \eqref{eq:pl-M-C-1}, and \eqref{eq:pl-M-C-2} that
\begin{equation}\label{eq:pl-M-C-3}
\frac{1-\beta}{2 \mu} \ge M \ge M-C=\frac{M(1-\beta-\mu M)-\frac{\beta^2 L \eta^2}{2}}{1-\beta-\mu M+\frac{\mu(\mu+L) \beta^2 \eta^2}{1-\beta}} \ge \frac{\frac{1-\beta}{16}(\gamma+\eta)}{\frac{9}{8}(1-\beta)}=\frac{\gamma+\eta}{18}>0.
\end{equation}

Next, we provide an upper bound on $\|\nabla f(x_k)\|^2$ using the $\mu$-PL condition \eqref{eq:pl-condition}. To this end, we note from the $L$-smoothness of $f$ and \eqref{eq:wk-xk} that
\begin{equation}\label{eq:pl-2}
\begin{aligned}
f(w_k) & \leq f(x_k)+\langle\nabla f(x_k), w_k-x_k\rangle+\frac{L}{2}\|w_k-x_k\|^2 \\
& =f(x_k)+\Big\langle\nabla f(x_k),-\frac{\beta \eta}{1-\beta} m_{k-1}\Big\rangle+\frac{L \beta^2 \eta^2}{2(1-\beta)^2}\|m_{k-1}\|^2 \\
& \leq f(x_k)+\frac{1}{2\mu}\|\nabla f(x_k)\|^2+\frac{\mu\beta^2\eta^2}{2 (1-\beta)^2}\|m_{k-1}\|^2+\frac{L \beta^2 \eta^2}{2(1-\beta)^2}\|m_{k-1}\|^2 \\
& =f(x_k)+\frac{1}{2\mu}\|\nabla f(x_k)\|^2+\frac{(\mu+L)\beta^2 \eta^2}{2(1-\beta)^2}\|m_{k-1}\|^2.
\end{aligned}
\end{equation}
Then it follows from the $\mu$-PL condition \eqref{eq:pl-condition} and \eqref{eq:pl-2} that
\[
\frac{1}{2 \mu}\|\nabla f(x_k)\|^2 \ge f(x_k)-f^* \ge f(w_k)-f^*-\frac{1}{2\mu}\|\nabla f(x_k)\|^2-\frac{(\mu+L)\beta^2 \eta^2}{2(1-\beta)^2}\|m_{k-1}\|^2,
\]
which further implies that
\begin{equation}\label{eq:pl-nabla_k}
\|\nabla f(x_k)\|^2 \ge \mu(f(w_k)-f^*)-\frac{\mu (\mu+L)\beta^2 \eta^2}{2(1-\beta)^2}\|m_{k-1}\|^2.
\end{equation}

Now we are ready to establish a recursive inequality for the Lyapunov functions $\{\varphi_k\}_{k\ge 1}$. Plugging $\rho=\frac{1-\beta}{2 L(\gamma+\eta)}$ into \eqref{eq:nc-3} gives
\[
\ebb[f(w_{k+1})] \leq \ebb[f(w_k)]+\frac{\beta^2 L \eta^2}{4(1-\beta)} \ebb[\|m_{k-1}\|^2]+\frac{L(\gamma+\eta)^2}{2} \sigma^2-M \ebb[\|\nabla f(x_k)\|^2],
\]
which together with the definition of $\{\varphi_k\}_{k\ge 1}$ \eqref{def:varphi} implies that
\begin{equation}\label{eq:pl-3}
\begin{aligned}
\ebb[\varphi_{k+1}]:=&\ebb\Big[f(w_{k+1})-f^*+C \sum_{j=1}^k \beta^{k-j}\|\nabla f(x_j)\|^2\Big]\\
=&\ebb[f(w_{k+1})-f^*]+C \ebb[\|\nabla f(x_k)\|^2]+\beta C \sum_{j=1}^{k-1} \beta^{k-1-j} \ebb[\|\nabla f(x_j)\|^2] \\
\leq & \ebb[f(w_k)-f^*]+\frac{\beta^2 L \eta^2}{4(1-\beta)} \ebb[\|m_{k-1}\|^2]+\frac{L(\gamma+\eta)^2}{2} \sigma^2-(M-C) \ebb[\|\nabla f(x_k)\|^2]\\
&+\beta C \sum_{j=1}^{k-1} \beta^{k-1-j} \ebb[\|\nabla f(x_j)\|^2].
\end{aligned}
\end{equation}
From $M-C>0$, applying \eqref{eq:pl-nabla_k} and Lemma \ref{lem:m-var} on \eqref{eq:pl-3}, we further obtain
\begin{equation}\label{eq:pl-4}
\begin{aligned}
\ebb[\varphi_{k+1}] \leq& [1-\mu(M-C)] \ebb[f(w_k)-f^*]+\Big(\frac{\beta^2 L \eta^2}{4(1-\beta)}+\frac{(M-C) \mu(\mu+L) \beta^2 \eta^2}{2(1-\beta)^2}\Big) \ebb[\|m_{k-1}\|^2] \\
&+\frac{L(\gamma+\eta)^2}{2} \sigma^2+\beta C \sum_{j=1}^{k-1} \beta^{k-1-j} \ebb[\|\nabla f(x_j)\|^2]\\
\leq & [1-\mu(M-C)] \ebb[f(w_k)-f^*]+\Big(\frac{\beta^2 L \eta^2}{2(1+\beta)}+\frac{(M-C) \mu(\mu+L) \beta^2 \eta^2}{1-\beta^2}+\frac{L(\gamma+ \eta)^2}{2}\Big) \sigma^2 \\
& +\Big(\beta C+\frac{\beta^2 L \eta^2}{2}+\frac{(M-C) \mu(\mu+L) \beta^2 \eta^2}{1-\beta}\Big) \sum_{j=1}^{k-1} \beta^{k-1-j} \ebb[\|\nabla f(x_j)\|^2].
\end{aligned}
\end{equation}
From the definition of $C$ \eqref{def:C}, it is straightforward to verify that
\begin{equation}\label{eq:pl-5}
\beta C+\frac{\beta^2 L \eta^2}{2}+\frac{(M-C) \mu(\mu+L) \beta^2 \eta^2}{1-\beta} = C(1-\mu M)\leq C[1-\mu(M-C)]
\end{equation}
Plugging \eqref{eq:pl-5} into \eqref{eq:pl-4} and noting \eqref{eq:pl-M-C-3} and $\mu\leq L$, we derive
\begin{equation}\label{eq:pl-6}
\begin{aligned}
\ebb[\varphi_{k+1}] & \leq[1-\mu(M-C)] \ebb[\varphi_k]+\Big(\frac{\beta^2 L \eta^2}{2(1+\beta)}+\frac{(M-C) \mu(\mu+L) \beta^2 \eta^2}{1-\beta^2}+\frac{L(\gamma+\eta)^2}{2}\Big) \sigma^2 \\
& \leq\Big(1-\frac{\mu(\gamma+\eta)}{18}\Big) \ebb[\varphi_k]+\Big(\frac{\beta^2 L \eta^2}{2(1+\beta)}+\frac{(\mu+L) \beta^2 \eta^2}{2(1+\beta)}+\frac{L(\gamma+\eta)^2}{2}\Big) \sigma^2 \\
& \leq\Big(1-\frac{\mu(\gamma+\eta)}{18}\Big) \ebb[\varphi_k]+\Big(\frac{3 \beta^2 L \eta^2}{2(1+\beta)}+\frac{L(\gamma+\eta)^2}{2}\Big) \sigma^2.
\end{aligned}
\end{equation}
Applying \eqref{eq:pl-6} recursively and noting that $\varphi_1=f(x_1)-f^*$ further leads to
\begin{equation*}
\begin{aligned}
&\ebb[f(w_{k+1})-f^*]\leq\ebb[\varphi_{k+1}]\\
\leq & \Big(1-\frac{\mu(\gamma+\eta)}{18}\Big)^k(f(x_1)-f^*)+\Big(\frac{3 \beta^2 L \eta^2}{2(1+\beta)}+\frac{L(\gamma+\eta)^2}{2}\Big) \sigma^2 \sum_{j=0}^{k-1}\Big(1-\frac{\mu(\gamma+\eta)}{18}\Big)^j \\
\leq & \Big(1-\frac{\mu(\gamma+\eta)}{18}\Big)^k(f(x_1)-f^*)+\Big(\frac{3 \beta^2 L \eta^2}{2(1+\beta)}+\frac{L(\gamma+\eta)^2}{2}\Big) \sigma^2 \cdot \frac{18}{\mu(\gamma+\eta)} \\
\leq & \Big(1-\frac{\mu(\gamma+\eta)}{18}\Big)^k(f(x_1)-f^*)+\Big(\frac{3 \beta^2 \eta}{1+\beta}+\gamma+\eta\Big) \frac{9L}{\mu} \sigma^2,
\end{aligned}
\end{equation*}
which completes the proof.
\end{proof}

\end{document}
